\subsection{Categorical reduction and the GRPO mean update}
\label{app:theory-grpo-mean}

Fix one prompt. Let the finite response categories be
$\mathcal A=\mathcal C\mathbin{\dot\cup}\mathcal N$, where the $m\ge2$
categories in $\mathcal C$ are correct execution modes and the $n\ge1$
categories in $\mathcal N$ are incorrect. Let
$p_a=\operatorname{softmax}(z)_a$, $P=\sum_{c\in\mathcal C}p_c$, and
$q_c=p_c/P$. Thus $q$ is exactly the conditional correct-mode distribution
$p_\theta^+(\cdot\mid x)$ of Section~\ref{sec:collapse}, and $P$ is the correct
mass $P_\theta^+(x)$.

\paragraph{Running example: three correct alternatives.}
\label{ex:theory-running}
Take three correct modes and one incorrect category, with
\[
 p=(0.30,0.15,0.05,0.50),\qquad P=0.50,\qquad q=(0.60,0.30,0.10).
\]
Half of all draws succeed, but among successful draws the three modes appear
in a $6{:}3{:}1$ ratio. Increasing $P$ and balancing $q$ are therefore
separate goals. We reuse this hypothetical distribution below to distinguish
the effects of the fresh objective, optimizer geometry, and replay.

Draw a group $A_1,\ldots,A_G\stackrel{\mathrm{iid}}{\sim}p$ with $G\ge2$,
write $R_i=\1[A_i\in\mathcal C]$, and let $M=\sum_iR_i$. Both GRPO variants
can be written, at the on-policy point where the PPO ratio equals one, as
\begin{equation}
  \widehat g_G
  =\frac1G\sum_{i=1}^G
   w(M)\left(R_i-\frac{M}{G}\right)\nabla_z\log p_{A_i}.
  \label{eq:group-estimator}
\end{equation}
For Dr.GRPO, $w(M)=1$ after absorbing the shared positive length scale. For
binary-reward GRPO, $w(M)$ is the reciprocal group reward standard deviation
on mixed groups. We absorb the same fixed response-length normalizer in both
variants; response-dependent row normalization is outside this reduction.
In either case $0<w(\ell)<\infty$ for $\ell\in\{1,\ldots,G-1\}$.
We define the entire update as zero on all-equal groups, without evaluating an
undefined reciprocal standard deviation; likewise, $w(M)M(G-M)$ below is
defined as zero when $M=0$ or $M=G$. We take explicit reference KL to be zero, as
in the experiments.

\begin{lemma}[A group baseline does not remove frequency amplification]
\label{lem:grpo-mean}
For $0<P<1$, the expected group update is
\begin{equation}
  \E[\widehat g_G]
  =c_G(P)\nabla_z P,
  \qquad
  c_G(P)=
  \frac{\E\!\big[w(M)M(G-M)\big]}{G^2P(1-P)}>0.
  \label{eq:expected-group-gradient}
\end{equation}
For Dr.GRPO, $c_G(P)=(G-1)/G$. Hence reward centering and binary reward
standardization change the speed, but not the expected direction, of the
binary-reward policy gradient.
\end{lemma}

\begin{proof}
Condition on the complete binary reward vector, whose sum is $M$. The score
identities give
\[
  \E[\nabla_z\log p_{A_i}\mid R_i=1]=\nabla_z\log P,
  \qquad
  \E[\nabla_z\log p_{A_i}\mid R_i=0]
  =\nabla_z\log(1-P).
\]
There are $M$ correct rows with centered reward $(G-M)/G$ and $G-M$
incorrect rows with centered reward $-M/G$. Therefore the conditional
expectation of the unaveraged sum in~\eqref{eq:group-estimator} is
\[
  w(M)\frac{M(G-M)}{G}
  \left(\nabla_z\log P-\nabla_z\log(1-P)\right)
  =w(M)\frac{M(G-M)}{GP(1-P)}\nabla_zP.
\]
The outer factor $1/G$ and expectation over
$M\sim\operatorname{Binomial}(G,P)$ yield~\eqref{eq:expected-group-gradient}.
For $w\equiv1$,
$\E[M(G-M)]=G(G-1)P(1-P)$, proving the Dr.GRPO value.
\end{proof}

The binomial identity also gives the denominator-free expression
\[
 c_G(P)=\frac{G-1}{G}\sum_{j=0}^{G-2}\binom{G-2}{j}
 w(j+1)P^j(1-P)^{G-2-j}.
\]
The binomial weights sum to one, so $c_G$ lies between
$(G-1)\min_j w(j+1)/G$ and $(G-1)\max_j w(j+1)/G$.
This establishes the smooth, positive extension to $P=0,1$ used in the
subsequent flow and convergence arguments.

The next lemma specializes the practical-estimator analysis of
\citet[Section~4.3 and Appendix~D, Theorem~5]{tajwar2026maxrl}
to our categorical notation; we include the derivation for completeness.

\begin{lemma}[Binary MaxRL has the same correctness-gradient direction]
\label{lem:maxrl-mean}
Let MaxRL assign
$A_i^{\mathrm{MaxRL}}=GR_i/M-1$ when $M>0$ and zero when $M=0$, as in
Section~\ref{sec:maxrl-method}. At the on-policy point, its expected update is
\begin{equation}
  \E[\widehat g_G^{\mathrm{MaxRL}}]
  =c_G^{\mathrm{MaxRL}}(P)\nabla_zP,
  \qquad
  c_G^{\mathrm{MaxRL}}(P)
  =\frac{1-(1-P)^{G-1}}{P}>0,
  \label{eq:maxrl-expected-gradient}
\end{equation}
with continuous endpoint values
$c_G^{\mathrm{MaxRL}}(0)=G-1$ and
$c_G^{\mathrm{MaxRL}}(1)=1$.
\end{lemma}

\begin{proof}
Condition on $M=m>0$. Each successful row has advantage $(G-m)/m$ and each
failed row has advantage $-1$. Using the same conditional score identities as
above, the expected unaveraged sum is
\[
  (G-m)\!\left(\nabla_z\log P-\nabla_z\log(1-P)\right)
  =\frac{G-m}{P(1-P)}\nabla_zP.
\]
After the outer factor $1/G$ and averaging over
$M\sim\operatorname{Binomial}(G,P)$,
\[
 \E[(G-M)\1\{M>0\}]
 =G(1-P)-G\Pr(M=0)
 =G\big[(1-P)-(1-P)^G\big].
\]
Substitution gives Equation~\eqref{eq:maxrl-expected-gradient}. Its endpoint
values follow by continuity (or the finite geometric-series identity), which
also gives the stated continuous extension at both boundaries.
\end{proof}

The same finite-series identity gives the potential used below:
\begin{equation}
 \Psi_{\mathrm{MaxRL}}(P)
 :=\int_0^P c_G^{\mathrm{MaxRL}}(s)\,ds
 =\sum_{k=1}^{G-1}\frac{1-(1-P)^k}{k}.
 \label{eq:maxrl-potential-finite-sum}
\end{equation}
In particular, $\Psi_{\mathrm{MaxRL}}(1)=\sum_{k=1}^{G-1}1/k$.
For binary rewards, best@\(k\) equals pass@\(k\), so this is a harmonic
mixture of sampling objectives for $k=1,\ldots,G-1$. The upper limit is
$G-1$, not $G$: the implemented centered estimator drops the entire update
on all-failure groups. Subtracting an unconditional score baseline even on
those groups instead gives the order-$G$ estimator. This distinction concerns
the exact finite-group expectation at the on-policy point; it does not assert
unbiasedness of subsequent clipped or adaptive-optimizer updates.

\paragraph{What the group normalization changes.}
For the running example with $G=4$,
$\nabla_zP=(0.15,0.075,0.025,-0.25)$. The two lemmas give
\[
 \begin{aligned}
 \mathbb E[\widehat g_4^{\rm Dr.GRPO}]&=0.75\,\nabla_zP,\\
 \mathbb E[\widehat g_4^{\rm MaxRL}]&=1.75\,\nabla_zP.
 \end{aligned}
\]
These are expectations over all possible groups, not updates from one
particular sample. MaxRL moves faster along the same direction at this
state, while the most common correct mode receives six times the logit
increment of the rarest. The next section asks what repeated unequal
increments do to the relative probabilities of equally rewarded modes.

